\subsection{An exact Gaussian construction for the path-sign comparison}
The graded coefficient objects of \fid{F43} and \fid{F45} admit an
explicit exponential-sum realization with the required opposite sign.
The same construction embeds point-Weil End/Hom data fully faithfully
in perverse sheaves. Its limitation is also exact: for this construction,
the pointwise estimate retains the entire coefficient-weight question.

\paragraph{Classical input and conventions.}
Let $k_0=\F_q$ have odd characteristic $p$, let $\ell\ne p$, and use a
characteristic-zero coefficient field $E$ containing the character
values. Let $\psi$ be a nontrivial additive character of $k_0$ and
$\psi_n=\psi\circ\Tr_{\F_{q^n}/k_0}$. Frobenius is geometric.
The classical inputs are the Artin--Schreier/Kummer trace construction
and its all-extension formula \cite[\emph{Sommes trig.},
1.7--1.9, especially (1.9.2)]{SGA45}, the rank-one Gaussian
cohomology in degree one \cite[Propositions~4.2--4.3,
pp.~196--197]{SGA45}, and its perfect pairing into $E(-1)$
\cite[Remark~4.4, p.~197]{SGA45}.
The original printed pages, not merely a modern transcription, were
inspected. We use the correct pairing/product identity in Remark~4.4,
not its misprinted absolute-value shorthand. No Gaussian size estimate
or Weil weight theorem is an input to the sign calculation below.

\begin{finding}[\fid{F50}: Gaussian path signs and the exact weight gate]
\label{f:gaussiansign}
There is a one-dimensional Gaussian Weil space $G_\psi$ such that
\[
 R\Gamma_c(\A^1_{\bar k_0},\mathcal L_\psi(x^2))=G_\psi[-1].
\]
For every effective finite-dimensional Weil space $C$ with invertible
Frobenius and algebraic eigenvalues, put
\[
 J_\psi(C)=\mathcal L_\psi(x^2)\otimes
                p^*(C\otimes G_\psi^\vee)[1],
 \qquad p:\A^1\longrightarrow\Spec k_0.
\]
Then $R\Gamma_c(\A^1,J_\psi(C))\simeq C$ naturally. Its trace
function reproduces $\Tr(F^n\mid C)$ exactly after summing over
$\A^1(\F_{q^n})$.

For an unbounded set of even path lengths $k$, its path
hypersurface complex is concentrated in degree $-1$. Its twisted
path sum is the negative trace of an explicit effective finite Weil
space, with zero error in \fid{F04}. Before the normalization in $J$,
the actual Gaussian path cohomology is concentrated in degree $k-1$.

The functor $J_\psi$ is exact and fully faithful into perverse
Weil sheaves in degree zero:
\[
 \Hom(J_\psi(C),J_\psi(C'))=\Hom_F(C,C').
\]
If $C$ is pure of weight $w$, then $J_\psi(C)$ is perverse pure
of weight $w$. Arithmetic Jordan blocks are allowed.

For any specified integer $w$, the pointwise bounds needed to
apply \fid{F03}/\fid{F17} to $J_\psi(C)$ and the reciprocal system
$J_{\psi^{-1}}(C^\vee(-1))$ hold under every embedding if and only if
every conjugate eigenvalue of $C$ has modulus $q^{w/2}$.
Thus the output-sign requirement is proved for this construction;
the two pointwise requirements retain precisely the original
coefficient-purity assertion.
\end{finding}

\begin{proof}
\par\smallskip\noindent\textbf{The Gaussian line without a weight estimate.}\quad
Write $\mathscr L_\psi=\mathcal L_\psi(x^2)$ and
$G_\psi=H_c^1(\A^1_{\bar k_0},\mathscr L_\psi)$.
The cited classical propositions give
\begin{equation}\label{eq:gaussline}
 R\Gamma_c(\A^1,\mathscr L_\psi)=G_\psi[-1],
 \qquad
 g_n:=\sum_{x\in\F_{q^n}}\psi_n(x^2)=-\gamma_\psi^n,
\end{equation}
where $\gamma_\psi$ is the Frobenius eigenvalue on the line.
They also give a perfect pairing
\begin{equation}\label{eq:gausspair}
 G_\psi\otimes G_{\psi^{-1}}\longrightarrow E(-1),
 \qquad \gamma_\psi\gamma_{\psi^{-1}}=q.
\end{equation}

The precise passage from the cited $\mathbb G_m$ calculation to this
$\A^1$ version is elementary. For the finite square map
$\pi:\A^1\to\A^1$, with $j:\mathbb G_m\hookrightarrow\A^1$,
\[
 \pi_*E=E\oplus j_!\mathcal L_\chi,
\]
where $\chi$ is the nontrivial quadratic character. The
anti-invariant summand has zero stalk at zero. The compact
cohomology of the nontrivial linear Artin--Schreier sheaf
$\mathcal L_\psi(x)$ is zero: it is a nontrivial character
summand of the Lang cover $\A^1\to\A^1$, whose source has
compact cohomology $E(-1)[-2]$ with trivial translation action.
Projection formula therefore identifies
$R\Gamma_c(\A^1,\mathscr L_\psi)$ with
$R\Gamma_c(\mathbb G_m,\mathcal L_\psi(x)\otimes\mathcal L_\chi)$.
Proposition~4.2 and Remark~4.4 apply.

The involution $x\mapsto-x$ acts by $-1$ on $G_\psi$,
since this is precisely the anti-invariant summand. For a
varying nonzero scalar $a$, the Gaussian line of
$\mathcal L_\psi(ax^2)$ is consequently $G_\psi$ tensored with
the quadratic Kummer line of $a$: use the etale change of
variable $x\mapsto\sqrt a\,x$ and its descent involution.
Completing the square also gives the exact Fourier identity
\[
 \sum_x\psi_n(x^2+xy)=g_n\psi_n(-y^2/4).
\]
This is a character-sum calculation, not positivity of the
Fourier operator.

\par\smallskip\noindent\textbf{Good even path lengths.}\quad
Let $A_k(h)$ have diagonal entries $1$, adjacent off-diagonal
entries $h/2$, and other entries zero. Then
\[
 x^{\mathsf T}A_k(h)x=\sum_i x_i^2+hP_k(x),\qquad
 P_k(x)=\sum_{i=1}^{k-1}x_ix_{i+1}.
\]
Set $d_k(h)=\det A_k(h)$, with $d_0=d_1=1$. Expansion gives
\begin{equation}\label{eq:gaussdetrec}
 d_k=d_{k-1}-\frac{h^2}{4}d_{k-2}.
\end{equation}
A root $h_0$ of $d_k$ is nonzero because $d_k(0)=1$.
At such a root $d_{k-1}(h_0)\ne0$: otherwise the recurrence
would make every preceding $d_j(h_0)$ zero, contradicting
$d_0=1$. Thus the singular matrix has rank exactly $k-1$
and its leading $(k-1)$-dimensional block is nonsingular.

For even $k=2m$, $d_k$ has degree $k$ and leading coefficient
\[
 \ell_k=(-1)^m/2^k.
\]
It is squarefree whenever $p\nmid k+1$.
Here is a proof retaining the characteristic hypothesis.
Define $U_0=1$, $U_1=2T$ and
$U_k=2TU_{k-1}-U_{k-2}$. The recurrence proves
\[
 d_k(h)=(h/2)^kU_k(1/h),\qquad
 U_k((z+z^{-1})/2)
   =\frac{z^{k+1}-z^{-(k+1)}}{z-z^{-1}}.
\]
If $p\nmid k+1$, the $2(k+1)$-st roots of unity are distinct.
Excluding $z=\pm1$ and identifying inverse pairs gives $k$
distinct roots of $U_k$; two pairs have the same sum
$z+z^{-1}$ only when they solve the same quadratic.
For even $k$ none has sum zero, because this would give
a fourth root of unity of order four while $k+1$ is odd.
Inverting the resulting nonzero roots proves squarefreeness.
Finally $U_{2m}(0)=(-1)^m$ gives the leading coefficient.

Call an even $k$ \emph{good} when $d_k$ is squarefree.
The sufficient condition supplies the unbounded good set
$k=2pr$, $r\ge1$. The condition is sufficient, not necessary:
$k=2$ is good in every odd characteristic.

\par\smallskip\noindent\textbf{Actual middle cohomology.}\quad
For a good even $k$, put $H_k=\{P_k=0\}\subset\A^k$.
We claim
\begin{equation}\label{eq:gaussmiddle}
 H_c^i(H_{k,\bar k_0},
       \mathcal L_\psi(\textstyle\sum_i x_i^2)|_{H_k})=0
       \quad\text{unless }i=k-1.
\end{equation}
This includes the singular vertex; it is not deduced from
the dimension of the hypersurface.

Consider $f:\A^1_h\times\A^k_x\to\A^1_h$ and
\[
 \mathscr Q=\mathcal L_\psi(\textstyle\sum_i x_i^2+hP_k(x)).
\]
Let $Z_k=\{d_k=0\}$, $U_k=\A^1\setminus Z_k$, and
write $j:U_k\hookrightarrow\A^1$, $i:Z_k\hookrightarrow\A^1$.
The finite scheme $Z_k$ is etale. Over $U_k$, etale-local
diagonalization and square roots of the unit coefficients,
together with \eqref{eq:gaussline} and Kunneth, show that the
only fiber compact cohomology is rank one in degree $k$.
Thus the generic system $T_k$ is lisse of rank one.

Near a root $h_0$, Schur complement splits off the nonsingular
leading block, leaving the scalar $d_k(h)/d_{k-1}(h)$.
At the root its free affine-line direction contributes
$E(-1)[-2]$. Its only fiber cohomology is therefore rank
one in degree $(k-1)+2=k+1$. Compact-support base change gives
\[
 R^kf_!\mathscr Q=j_!T_k,\qquad
 R^{k+1}f_!\mathscr Q=i_*V_k,
\]
with all other direct-image sheaves zero and $V_k$ rank one
on $Z_k$. The first equality follows from its zero stalks
at $Z_k$. The compact-support formalism used here follows
from proper base change and compactification
\cite[Theorem~17.10, p.~116; Proposition~18.4,
pp.~121--122]{Milne}; Kunneth and projection formula are
\cite[Theorem~22.1, Remarks~22.3(b),(d), and
Proposition~22.5, pp.~135--137]{Milne}.

At every simple root, the Schur-complement scalar has order
one. Its Gaussian factor has quadratic Kummer monodromy
$-1$, while the nonsingular block is locally unramified.
Hence $T_k$ is geometrically nonconstant. On the connected
nonproper curve $U_k$,
\[
 H_c^0(U_k,T_k)=H_c^2(U_k,T_k)=0.
\]
The first assertion concerns proper-support sections of a
lisse sheaf; the second follows by duality from the absence
of geometric invariants of the dual line
\cite[\emph{Sommes trig.}, 1.18(b)--(d),
pp.~178--179]{SGA45}. No weight bound enters.
The compact-support Leray spectral sequence has only the
terms $(p,q)=(1,k)$ and $(0,k+1)$, both in total degree
$k+1$. Thus $R\Gamma_c(\A^1\times\A^k,\mathscr Q)$ is
concentrated in degree $k+1$.

Integrating $h$ first gives zero when $P_k(x)\ne0$, by the
linear Artin--Schreier calculation, and $E(-1)[-2]$ when
$P_k(x)=0$. Base change and projection formula give
\[
 R\Gamma_c(\A^1\times\A^k,\mathscr Q)
 \simeq
 R\Gamma_c(H_k,\mathcal L_\psi(\textstyle\sum_i x_i^2))(-1)[-2].
\]
This proves \eqref{eq:gaussmiddle}.

\par\smallskip\noindent\textbf{An explicit effective output spectrum.}\quad
Let $C_k$ be the smooth projective completion of
$y^2=d_k(h)$. It is geometrically connected.
The double cover has $k$ simple branch points and no
branch at infinity, so Riemann--Hurwitz gives genus
$(k-2)/2$. Let $\eta_k$ be the quadratic Weil line whose
extension trace is $\chi_n(\ell_k)$; it is the anti-invariant
line of the two points at infinity. Put
\[
 W_k=H^0(Z_{k,\bar k_0},
             \mathcal L_\chi(d_{k-1}|_{Z_k})).
\]
This is well defined by \eqref{eq:gaussdetrec} and has
dimension $k$. It retains the actual Frobenius permutation
of the roots, which need not be rational over $k_0$.
Define the effective Weil space
\begin{align}
 B_{k,\psi}={}&G_\psi^{\otimes k}\otimes
                 (H^1(C_k,E)\oplus\eta_k)\notag\\
 &\oplus E(-1)\otimes G_\psi^{\otimes(k-1)}\otimes W_k .
 \label{eq:gaussoutput}
\end{align}
Its dimension is $2k-1$ and its eigenvalues have genuine
positive multiplicities.

For $Q=q^n$ and
$N_{k,n}^\psi=\sum_{P_k=0}\psi_n(\sum_i x_i^2)$,
orthogonality and diagonalization give
\begin{align}
 QN_{k,n}^\psi
 ={}&g_n^k\sum_{h\in\F_Q}\chi_n(d_k(h))\notag\\
 &+Qg_n^{k-1}
       \sum_{\substack{h\in\F_Q\\d_k(h)=0}}
                         \chi_n(d_{k-1}(h)).
 \label{eq:gausspenciltrace}
\end{align}
Indeed a nonsingular quadratic form contributes its
determinant character times $g_n^k$; at a root Schur
complement gives one free affine-line factor and the
displayed nonsingular minor. The projective curve trace
formula gives, without a curve weight estimate,
\[
 \sum_h\chi_n(d_k(h))
 =-\Tr(F^n\mid H^1(C_k))-\chi_n(\ell_k).
\]
Its affine point count is $Q$ plus the sum on the left;
its infinity fiber has $1+\chi_n(\ell_k)$ points, while
$H^0=E$ and $H^2=E(-1)$.
Since $k$ is even, $g_n^k=\gamma_\psi^{kn}$ and
$g_n^{k-1}=-\gamma_\psi^{(k-1)n}$. Therefore
\begin{equation}\label{eq:gaussnegativesign}
 QN_{k,n}^\psi=-\Tr(F^n\mid B_{k,\psi})
                      \qquad(n\ge1).
\end{equation}
This identifies the character of the twisted middle
cohomology with \eqref{eq:gaussoutput}; no canonical
splitting or arithmetic semisimplicity is asserted.

The algebraicity and embedding requirements are explicit.
The traces of $H^1(C_k)$ are the rational integers
$1+q^n-\#C_k(\F_{q^n})$. Newton identities in its finite
dimension give a rational characteristic polynomial, hence
an algebraic embedding-stable spectrum. Frobenius on $W_k$
has finite order, $\eta_k$ is quadratic, and
$\gamma_\psi=-g_1$ is cyclotomic. Conjugating
\eqref{eq:gaussoutput} therefore gives exactly the
construction for the conjugated character $\psi$.
No Weil integrality or weight theorem is needed for this step.

At $k=2$, the union of the two axes gives
$N_{2,n}^\psi=2g_n-1$, so
\[
 QN_{2,n}^\psi=-2Q\gamma_\psi^n-Q.
\]
This also checks every sign in \eqref{eq:gaussoutput};
the relation $\gamma_\psi^2\chi(-1)=q$ follows from
\eqref{eq:gausspair} and quadratic scaling.
The good-length condition matters. In characteristic five,
$k=4$ has $d_4=(h^2-1)^2$ and $d_3(\pm1)=3$.
For $q=5$, \eqref{eq:gausspenciltrace} gives
\[
 QN_{4,n}^\psi=5^{3n}-2\cdot5^{2n}
                         -2(-5\gamma_\psi^3)^n.
\]
Here $\gamma_\psi^2=5$. The root $125$ has positive
coefficient and strictly largest modulus, so this bad
length cannot have an exact negative effective spectrum.
The proof supplies an unbounded good even set, not all
even lengths in every characteristic.

\par\smallskip\noindent\textbf{The size follows from the sign theorem.}\quad
For $t_n(x)=\psi_n(x^2)$ the pointwise bound is $1$.
Thus $a=0,d=1$ in \fid{F03}.
Equation~\eqref{eq:gaussline} is its negative rank-one
input spectrum, and \eqref{eq:gaussnegativesign} supplies
the opposite output on the unbounded good set, with zero
error. Finding~\fid{F03} gives $|\gamma_\psi|\le q^{1/2}$.
Apply the same argument to $\psi^{-1}$ and then
\eqref{eq:gausspair}: both absolute values equal $q^{1/2}$.
Every embedding has its corresponding nontrivial character
and output spectrum, so this holds for every complex
absolute value. This recovers the classical Gaussian
result as an application test of the sign criterion.

\par\smallskip\noindent\textbf{Normalization and the exact target trace.}\quad
For an effective finite Weil space $C$, projection formula
and evaluation $G_\psi\otimes G_\psi^\vee\to E$ give
\[
 R\Gamma_c(\A^1,J_\psi(C))\simeq C.
\]
The shifts $[-1]$ and $[1]$ cancel. Its trace function is
\begin{equation}\label{eq:gaussnormalizedtrace}
 t_{C,n}(x)=-\gamma_\psi^{-n}
            \Tr(F^n\mid C)\psi_n(x^2),
 \qquad M_{C,n}=\Tr(F^n\mid C).
\end{equation}
For a good even $k$,
\begin{align*}
 R\Gamma_c(H_k,J_\psi(C)^{\boxtimes k}|_{H_k})
 \simeq{}&C^{\otimes k}\otimes G_\psi^{\vee\otimes k}\\
 &\otimes
 R\Gamma_c(H_k,\mathcal L_\psi(\textstyle\sum_i x_i^2))[k].
\end{align*}
Its only degree is $(k-1)-k=-1$. Thus its effective
twisted output is
\[
 B_{k,C}=C^{\otimes k}\otimes
                    G_\psi^{\vee\otimes k}\otimes B_{k,\psi},
 \qquad
 QN_{k,C,n}=-\Tr(F^n\mid B_{k,C}).
\]
This proves the exact opposite sign with no boundary error.
It concerns effective coefficient spaces, including the
even-degree total realization of \fid{F45}; it does not
declare tensor powers of arbitrary signed virtual complexes
to be effective.

\par\smallskip\noindent\textbf{The pointwise and dual gates.}\quad
For nonzero $C$ and a fixed embedding put
\[
 b_{\max}=2\log_q\max_\alpha|\alpha|,\qquad
 b_{\min}=2\log_q\min_\alpha|\alpha|,
 \qquad \alpha\in\operatorname{Spec}(F\mid C).
\]
These are real thresholds, not yet asserted integer weights.
The Gaussian result just proved gives
\[
 \sup_x|t_{C,n}(x)|
       =q^{-n/2}|\Tr(F^n\mid C)|.
\]
There is a constant $B$ independent of $n$ with
$|t_{C,n}(x)|\le Bq^{na/2}$ if and only if
\begin{equation}\label{eq:gausspointwisegate}
 \max_\alpha|\alpha|\le q^{(a+1)/2}.
\end{equation}
Sufficiency uses $B=\dim C$; necessity is \fid{F02}.
Jordan blocks do not introduce polynomial factors into
traces of powers. The sharp real exponent is therefore
$b_{\max}-1$.

The reciprocal system is
$J_{\psi^{-1}}(C^\vee(-1))$. Its total spectrum is
$\{q/\alpha\}$, with the required perfect pairing into
$E(-1)$, and its sharp real pointwise exponent is
$1-b_{\min}$. Relative to the formally matching exponent
$1-b_{\max}$, the exact gap is $b_{\max}-b_{\min}$.

To invoke \fid{F03}/\fid{F17} as actually stated, fix one
integer target weight $w$ and set $a=w-1$.
The two pointwise bounds at exponents $w-1$ and $1-w$
hold under every embedding precisely when
\begin{equation}\label{eq:gausscommonweight}
 b_{\max,\sigma}\le w\le b_{\min,\sigma}
                    \quad\text{for every }\sigma.
\end{equation}
Equivalently all conjugate eigenvalues have modulus
$q^{w/2}$. Width zero separately in every embedding is
insufficient if the radii differ between embeddings:
the same integer $w$ must work throughout.

\par\smallskip\noindent\textbf{Hom and conditional perverse purity.}\quad
For $C,C'$, the equal shifts and the same invertible
rank-one sheaf and Gaussian line cancel from
$\Hom(J_\psi(C),J_\psi(C'))$. A map between constant
sheaves on connected geometric $\A^1$ is exactly a
linear map of their fibers. Weil compatibility is exactly
Frobenius commutation. Therefore, naturally and
compatibly with composition,
\[
 \Hom(J_\psi(C),J_\psi(C'))=\Hom_F(C,C').
\]
This is both perverse-heart Hom and degree-zero derived
Weil Hom; no higher-Ext comparison is claimed.
The construction is exact. If $C$ is pure of weight $w$,
its unshifted lisse coefficient has weight $w-1$ and its
shift by $[1]$ has complex weight $w$ in BBD's convention
\cite[5.1.8, pp.~126--127]{BBD}.
It is perverse because the base is a smooth curve.
Arithmetic Jordan blocks remain allowed; geometrically
the constituents are copies of one rank-one local system.
\end{proof}

\paragraph{Two precise modes of retaining the actual actions.}
At a fixed arithmetic specialization of \fid{F46}, including
the collision parameter, its actual $H,K$ give
\[
 \End(JH)=A,\qquad \Hom(JH,JK)=M,\qquad
 \End(JK)=\End_A(M).
\]
Precomposition and postcomposition retain the specified
right generators and ordinary left operators, including
the arithmetic collision nilpotent. With \fid{F46}'s
chosen pure weight-zero Weil data, $JH$ and $JK$ are
actual pure weight-zero perverse Weil sheaves.
Their Frobenius powers are those of the fixed operators;
they are not extension-by-extension quantum-parameter
substitutions.

The sign statement applies to identity traces of effective
objects. An arbitrary inserted trace $\Tr(TF^n\mid C)$ is
different: its Frobenius coefficients need not be effective
multiplicities, and its tensor-inserted output need not have
the required sign. Retaining $T$ as a morphism does not
establish that extra assertion.

For \fid{F43}'s graded model, take a point-Weil line $U$
with eigenvalue $q^{1/2}$ and replace every degree-$e$
edge by $U^e$. Each named kernel then has trace matrix
equal to its actual Laurent matrix evaluated at
$v=q^{n/2}$. Odd degrees use a half-weight arithmetic
Weil line; they are not asserted to be integral Tate
motives. Even degrees use the usual Tate lines.
For any constant coefficient space $V$,
\[
 J_\psi(C\otimes V)\simeq p^*V\otimes J_\psi(C).
\]
Thus the kernel functors and their explicitly ordered
word compositions transport through $J$. All the actual
graded $K_0$ matrices of $P_1,P_2$ and the stated left
operators remain those of \fid{F43}, with its opposite
row-composition convention for the left action.
This does not choose a coherent lift of every algebra
relation or commutation beyond that established there.
The word coefficient objects change with the word.

There is no strong monoidal claim:
$J(E)$ is not the tensor unit, and
\[
 J(C)\otimes J(C')
 =\mathcal L_{\psi^2}(x^2)\otimes
     p^*(C\otimes C'\otimes G_\psi^{\vee\otimes2})[2],
\]
whereas $J(C\otimes C')$ has one Gaussian factor and
shift $[1]$. Nor have distinct geometric simple labels
or a designated canonical basis been constructed.
These exact scope restrictions agree with \fid{F42}
and \fid{F48}.

\paragraph{The actual word objects and their dual gap.}
For a nonempty right word of length $L$ and cyclic
transition count $2s$, write
\[
 T_w(v)=20D^L+16D^{L-s}
          +4D^{L-2s}(D-2)^{2s}
       =\sum_{r=-L}^{L}c_rv^{2r},\qquad c_r\ge0.
\]
The total effective Weil space of \fid{F45} is
$C_w=\bigoplus_r E(-r)^{c_r}$.
Equation~\eqref{eq:gaussnormalizedtrace} reproduces
$T_w(q^{n/2})$ exactly, and its path output has the
proved negative effective spectrum with zero error.
Both extreme coefficients are nonzero, already from
$20D^L$. Hence $b_{\max}=2L$, $b_{\min}=-2L$.
The sharp pointwise exponents are $2L-1$ for the original
system and $2L+1$ for its reciprocal system. At target
weight $2L$, \fid{F17} would need the latter exponent
$1-2L$: the exact gap is $4L$.

For each homogeneous piece $E(-r)^{c_r}$ the two exponents
are instead $2r-1$ and $1-2r$. The criterion recovers
exactly its known pure weight $2r$. The analogous statement
holds for every homogeneous path coefficient of \fid{F43},
using $U^e$ when its degree is odd. Its actual finite path
labels and explicit-word bounds are available from that
finding. No new all-rank positive rule, canonical-basis
identification or binary-input PH1 consequence follows.

\paragraph{The direct positive graph-count route has a different sign.}
Putting the effective coefficient trace $S_n=\Tr(F^n\mid C_w)$
at the origin and zero elsewhere on $\A^1$ gives
\[
 M_n=S_n,\qquad QN_{k,n}=QS_n^k.
\]
This nonzero positive-multiplicity Tate spectrum cannot be
the negative of an effective spectrum for every extension:
their sum would contradict \fid{F02}. This excludes the
exact sign for that point-supported construction. Its
small support still permits loose error estimates with
empty output; no broader \fid{F04} obstruction is inferred
from this example alone.

Affine padding gives a sharper quantitative failure.
For a base $\A^d$ with $d\ge1$, take the constant trace
$t_n(x)=Q^{-d}S_n$, so $M_n=S_n$ and the sharp pointwise
exponent is $a=2L-2d$. For even $k=2m$, separating the
odd and even coordinates makes $P_k$ a nondegenerate
bilinear pairing between two $md$-dimensional spaces:
the block coefficient matrix is triangular with invertible
diagonal. The zero first vector contributes $Q^{md}$
choices, and every nonzero first vector imposes one
nonzero linear equation. Hence
\[
 \#H_k(\F_Q)=Q^{kd-1}+Q^{kd/2}-Q^{kd/2-1},
\]
and
\[
 QN_{k,n}=S_n^k
             \bigl(1+Q^{1-kd/2}-Q^{-kd/2}\bigr).
\]
For $kd>2$ its highest-radius root is $q^{kL}$ with
positive coefficient $c_L^k$; for $kd=2$ that coefficient
is twice as large. An added effective spectrum cannot
cancel it. If this were minus an effective spectrum plus
an \fid{F04} error, \fid{F02} would therefore force
\[
 q^{kL}\le q^{k(a+d)/2+\eta_k}
              =q^{k(L-d/2)+\eta_k},
 \qquad \eta_k\ge kd/2.
\]
The cost is linear in $k$, incompatible with the required
sublinear overhead. Full extension-degree progressions
do not remove the positive top Tate root.
Increasing $a$ by $d$ can absorb the cost, but weakens
the resulting weight estimate by $d$. This is the exact
boundary for the specified normalized affine-padding route.

% Separate paragraph for Part III.
\paragraph{\fid{F50} in Grothendieck's language.}
The concrete schemes are an affine line, the path quadric,
its quadratic pencil, a finite etale degeneracy divisor,
and a smooth projective double cover. Artin--Schreier
and Kummer torsors supply the actual coefficient sheaves.
Compact-support base change, Kunneth, duality and Leray
prove the single-degree statement; the trace formula and
Tate twist give the opposite sign and factor $q^n$.
The exact primary inputs are SGA~$4\frac12$,
\emph{Sommes trig.}, 1.9.2, 1.18 and 4.2--4.4, and
the Milne locators displayed above; the trace framework
also agrees with SGA~5, III, 4.7 \cite{SGA45,Milne,SGA5}.
The fixed-parameter action language is that of equivariant
modules, SGA~3, I, 4.7.2--3, and cote~90, scan~17
\cite{SGA3,Archive}.

The normalized construction gives a fully faithful
degree-zero perverse realization of the actual
\fid{F46} End/Hom data, while remaining a module functor
for constant coefficients rather than a strong monoidal
functor. It retains arithmetic nilpotents and does not
create distinct geometric simple labels. Its graded
version retains the actual \fid{F43} coefficient kernels
with the stated $K_0$ and coherence qualifications.
The opposite-sign comparison is proved; the remaining
pointwise and dual requirements are exactly
\eqref{eq:gausspointwisegate}--\eqref{eq:gausscommonweight}.
A further advance would require replacing the constant
unknown coefficient factor by a varying local system
with an independently established pointwise bound and
the required actual operation comparisons.

% Bibliography expansion (existing keys, for parent integration):
% SGA45: add Applications de la formule des traces aux sommes
% trigonometriques, 1.7--1.9, 1.18, Props4.2--4.3 pp196--197,
% Remark4.4 p197; original author scan
% https://publications.ias.edu/sites/default/files/Number32.pdf.
% Milne: additionally Theorem17.10 p116, Proposition18.4 pp121--122,
% Theorem22.1, Remarks22.3(b),(d), Proposition22.5 pp135--137.
